% 
%  chapter2.tex
%  ReportISEL
%  
%  Created by Matilde Pós-de-Mina Pato on 2012/10/09.
%  Modified by José Saldanha, Ricardo Pinto and Afonso Jesus on 2025/06/02
%
\chapter{Exercício 2}
\label{cha:exercise_2}

\section{Implementação da comunicação}

Pretende-se implementar o envio e receção de uma mensagem entre \textbf{Alice} e \textbf{Bob}, usando um
esquema de \textbf{cifra/decifra assimétrica} e a função de hash \texttt{SHA256}.
A chave pública de Bob (\texttt{KeBob}) é usada por Alice para cifrar a mensagem, e a chave privada de Bob (\texttt{KdBob}) é usada para a decifrar.

\subsection*{Envio (Alice)}

\begin{verbatim}
# Cifra a mensagem com a chave pública de Bob
openssl rsautl -encrypt -inkey KeBob.pem -pubin -in m -out c

# Calcula o hash SHA256 da mensagem original
openssl dgst -sha256 -binary m > h

# Concatena ambos os ficheiros
cat c h > ch
\end{verbatim}

Assim, o ficheiro \texttt{ch} contém o criptograma e o hash:
\[
ch = E_a(Ke_{Bob})(m) || SHA256(m)
\]

\subsection*{Receção (Bob)}

\begin{verbatim}
# Divide o ficheiro recebido
head -c <tam_c> ch > c
tail -c 32 ch > h

# Decifra com a chave privada de Bob
openssl rsautl -decrypt -inkey KdBob.pem -in c -out m

# Recalcula o hash da mensagem decifrada
openssl dgst -sha256 -binary m > h2

# Compara os hashes
cmp h h2 && echo "Mensagem íntegra" || echo "Falha de integridade"
\end{verbatim} 

No final do processo, Bob obtém a mensagem original $m$ e verifica se o hash corresponde ao valor enviado.

\section{Análise de confidencialidade e integridade}

\subsection*{Confidencialidade}

A cifra assimétrica garante \textbf{confidencialidade}, pois a mensagem é cifrada com a chave pública de Bob:
\[
m = D_a(Kd_{Bob})(c)
\]
Apenas Bob, possuindo a chave privada correspondente, consegue decifrar a mensagem. Um atacante que intercepte o criptograma \(c\) não consegue recuperar \(m\) sem conhecer \(Kd_{Bob}\).

\subsection*{Integridade e autenticidade}

O hash SHA256 permite verificar se a mensagem foi alterada após a transmissão. Contudo, este esquema \textbf{não garante integridade de forma robusta} nem autenticidade de Alice:

\begin{itemize}
  \item O hash é enviado em claro juntamente com o criptograma:
  \[
  ch = E_a(Ke_{Bob})(m) \;||\; SHA256(m)
  \]
  Um atacante poderia substituir o criptograma \(c\) e criar um novo par \((c', h')\), o que poderia induzir Bob a aceitar uma mensagem alterada.
  
  \item Para garantir \textbf{integridade} e \textbf{autenticidade}, recomenda-se que:
  \begin{itemize}
    \item o hash seja protegido com uma chave secreta (por exemplo, usando HMAC), ou
    \item a mensagem seja assinada digitalmente por Alice:
    \[
    s = S(Kd_{Alice})(SHA256(m))
    \]
    permitindo que Bob verifique a assinatura:
    \[
    V(Ke_{Alice})(s, SHA256(m)) = \text{true}
    \]
  \end{itemize}
\end{itemize}